\documentclass[10pt,stdletter,dateno]{newlfm}
\usepackage{kpfonts}
\usepackage{url}

\makeatletter
\g@addto@macro{\ps@ltrhead}{%
	\renewcommand{\headrulewidth}{0pt}%
	\renewcommand{\footrulewidth}{0pt}%
}
\g@addto@macro{\ps@othhead}{%
	\renewcommand{\headrulewidth}{0pt}%
	\renewcommand{\footrulewidth}{0pt}%
} 
\makeatother


\widowpenalty=1000
\clubpenalty=1000

\usepackage{fontspec}
% use sans family font
\renewcommand{\familydefault}{\sfdefault}

%\newlfmP{headermarginskip=1pt}
\newlfmP{sigsize=5pt}
\newlfmP{dateskipafter=10pt}
\newlfmP{addrfromphone}
\newlfmP{addrfromemail}
\PhrPhone{Phone}
\PhrEmail{Email}

\namefrom{Sajed Zarinpour Nashroudkoli}
\addrfrom{%
	\today\\[3pt]
	30 Dashti Avenue\\
	Komeil Street\\
	Tehran 1353673335\\
	IRAN
}

\phonefrom{+98-935-218-82-08}
\emailfrom{sajed.zarinpour@gmail.com}

\addrto{%
	Karlsruher Institut für Technologie (KIT)\\
	Fakultät für Mathematik\\
	Institut für Angewandte und Numerische Mathematik\\
	Arbeitsgruppe 3: Wissenschaftliches Rechnen\\
	Englerstr. 2\\
	D-76133 Karlsruhe\\
	GERMANY
}

\greetto{
	Subject: Motivation Letter for Doctoral Researcher (f/m/d) Project B7 “Dynamics of electrical depolarization waves in the heart”\\\vspace{1em}
	Dear Prof. Dr. Christian Wieners and the Selection Committee,
}
\closeline{Sincerely,}
\begin{document}
	\begin{newlfm}
		
		I am writing to express my strong interest in the Doctoral Researcher position for Project B7, "Dynamics of electrical depolarization waves in the heart," at CRC 1173. With a Master's degree in Numerical Analysis from the Institute for Advanced Studies in Basic Sciences (IASBS), Iran, I have developed a strong foundation in combining machine learning methodologies with traditional numerical techniques for solving partial differential equations (PDEs).  The opportunity to contribute my expertise in numerical modeling and programming to the interdisciplinary challenges of simulating heart function within Project B7 is particularly compelling to me.
		
%		During my Master's program, my research was initially motivated by the challenges of biophysical modeling, specifically in simulating breast deformation under compression. While my thesis ultimately focused on the mathematical methodology – developing a novel neural network approach to solve PDEs with uncertainties – this initial project, along with explorations into melanoma cancer and cellular growth modeling, ignited my passion for applying numerical techniques to biological systems. This experience provided me not only with a deepened understanding of numerical methods, finite element methods, and computational modeling, directly applicable to the complexities of depolarization waves and heart muscle contraction, but also honed my programming skills and creative problem-solving abilities, exemplified by the development of a custom periodic padding layer for TensorFlow/Keras to efficiently handle domain uncertainty in PDEs.  Furthermore, presenting my research findings in English at the IASBS Student Presentations in English further solidified my communication skills in an academic setting.

		During my Master’s studies, I worked extensively on numerical methods for solving PDEs with uncertainty, utilizing deep learning techniques. Initially motivated by challenges in biophysical modeling, I explored applications in breast deformation under compression, melanoma cancer modeling, and cellular growth simulations. While my thesis ultimately focused on developing a novel neural network approach for solving PDEs with domain uncertainties, these projects provided me with valuable experience in computational modeling and finite element methods—both of which are directly applicable to simulating depolarization waves and heart muscle contraction. Additionally, I implemented a custom periodic padding layer in TensorFlow/Keras to efficiently handle domain uncertainties, demonstrating my ability to develop innovative numerical techniques. Furthermore, presenting my research findings at the IASBS Student Presentations in English further solidified my communication skills in an academic setting. 
		
		Beyond my academic pursuits, my experience as a software developer helped me cultivate valuable teamwork, problem-solving, and time management skills.  Collaborating within multidisciplinary teams to deliver projects under tight deadlines has honed my communication and adaptability. This period of professional software development, undertaken alongside personal responsibilities including caring for my elderly father, broadened my skillset and reinforced my dedication to focused, impactful work.  With my family situation now stabilized and with renewed energy, I am eager to fully dedicate myself to research and return to my academic passion, specifically within a field that combines mathematical rigor with real-world impact, such as Project B7.
		
%		I am particularly drawn to this position due to its focus on the exciting intersection of mathematics, biomedical engineering, and advanced computation.   My engagement to this field extends beyond my formal education. As evidenced by my participation in various seminars and conferences during recent years, I have actively followed advancements in biophysical modeling and numerical methods. This continuous self-driven learning has further solidified my desire to contribute to cutting-edge research in this area, making Project B7 and CRC 1173 an ideal next step for me. The opportunity to contribute to the modeling and numerical simulation of cardiac dynamics within the inspiring and interdisciplinary environment at KIT aligns perfectly with my long-term aspiration to conduct impactful research at the interface of mathematics and healthcare.
		
		
		What particularly draws me to Project B7 is its focus on the mathematical and computational challenges associated with cardiac dynamics. My participation in international conferences and schools on inverse problems, biophysical modeling, and numerical analysis has further solidified my interest in this field. I am excited about the prospect of contributing to developing efficient numerical methods for simulating heart function, particularly in a collaborative research environment such as CRC 1173, which aligns perfectly with my long-term aspiration to conduct impactful research at the interface of mathematics and healthcare.
		
		With my background in numerical analysis, robust programming expertise, and deep-seated passion for applying mathematics to understand biological systems, I am confident in my ability to make significant contributions to Project B7. 
%		I am eager to immerse myself in the stimulating research environment at CRC 1173, learn from the pioneering work being conducted, and further develop as a researcher in this dynamic field.
		
		Thank you sincerely for considering my application. I look forward to the opportunity to join your research group, collaborate with esteemed researchers, and further develop my expertise in computational mathematics and biomedical applications.  
%		I am very enthusiastic about the possibility of joining your research group and would welcome the opportunity to discuss my application further.
		
		\vspace{1em}
		P.S. I am in the process of obtaining my official Master's degree certificate, which I anticipate receiving within the next 3-4 months. Unofficial transcripts and confirmation from my supervisor, Dr. Khadijeh Nedaiasl, are available immediately and I will gladly provide the official certificate as soon as it is issued.
		
	\end{newlfm}
\end{document}